\documentclass[11pt]{article}

% ARR / ACL review mode.  Keep the official template unchanged.
\usepackage[review]{acl}

\usepackage{times}
\usepackage{latexsym}
\usepackage[T1]{fontenc}
\usepackage[utf8]{inputenc}
\usepackage{microtype}
\usepackage{inconsolata}
\usepackage{graphicx}
\usepackage{amsmath}
\usepackage{amssymb}
\usepackage{booktabs}
\usepackage{multirow}
\usepackage{xcolor}
\usepackage{algorithm}
\usepackage{algorithmic}
\usepackage{subcaption}
\usepackage{enumitem}
\usepackage{array}
\usepackage{makecell}
\usepackage{xurl}
\usepackage{placeins}

\graphicspath{{../../experiments/main_experiment/paper_assets/figures/}}

% Keep float-only appendix pages top-aligned instead of vertically centered.
\makeatletter
\setlength{\@fptop}{0pt}
\setlength{\@fpsep}{8pt plus 2pt}
\setlength{\@fpbot}{0pt plus 1fil}
\setlength{\@dblfptop}{0pt}
\setlength{\@dblfpsep}{8pt plus 2pt}
\setlength{\@dblfpbot}{0pt plus 1fil}
\makeatother

\newcommand{\samcgs}{SA-MCGS}
\newcommand{\naive}{Naive}
\newcommand{\riskany}{Risk-any}
\newcommand{\riskall}{Risk-all}
\newcommand{\rootthree}{Root@3}
\newcommand{\scc}{SCC}
\newcommand{\oc}{OC}

% Revision view: earlier additions blue/green, October 1 changes red; original deletions struck.
\input{revision_markup}

\title{Risk Subgraph Extraction in Cyclic Document Graphs\\
with Structure-Aware Monte Carlo Search}

\author{Anonymous ARR Submission}

\begin{document}
\maketitle

\begin{abstract}
Long documents such as contracts, statutes, software package metadata, and corporate filings often contain risks that emerge through mutually dependent records.  LLM-guided Monte Carlo tree search (MCTS) can improve structured document reasoning, but standard tree rollouts become unstable on strongly connected components (SCCs), where the same logical record may be revisited indefinitely.  We propose \samcgs{}, a structure-aware Monte Carlo graph search method that localizes SCCs, performs local LLM rollouts with relation-first evidence memory, and collapses each cyclic region into an inspectable risk subgraph.  We evaluate a locked critical structural benchmark covering four authority-backed domains: Debian dependencies, SEC EX-21 subsidiary disclosures, German Civil Code cross-references, and CUAD contracts.  \revredreplace{Across 320 model-case pairs and four LLMs, \samcgs{} improves \rootthree{} from 58\% to 81\%, \riskany{} from 79\% to 98\%, and \riskall{} from 62\% to 78\% over an oracle-risk Top-3 full-\scc{} \naive{} baseline selected from ten attempts.}{Across five LLMs and 400 model-cases, \samcgs{} raises mean \rootthree{} and \riskall{} over local evidence aggregation (LEA) by 6.75 and 15.25 percentage points, respectively. The corresponding gains over one-shot prompting are 46.00 and 38.75 points. Returned subgraphs remove 55.70\% of SCC nodes on average.} \revredreplace{The gains are larger under long-context pressure, and the returned subgraphs retain core evidence while removing more than half of the original SCC nodes.}{Component ablations isolate the effect of core replacement, while cost measurements and a blinded expert audit assess computational efficiency and the quality of the extracted evidence.}
\end{abstract}

\section{Introduction}
\label{sec:intro}

Many high-stakes document analysis tasks are graph problems hidden inside text \citep{hendrycks2021cuad,koreeda2021contractnli,boulet2021measuring}.  In contracts, one clause may define a deadline, another may suspend it, and a later amendment may restore only part of the obligation.  Similar dependencies appear when statutory exceptions are narrowed by distant cross-references \citep{boulet2021measuring,holzenberger2020sara}, or when software packages preserve old interfaces while dependent packages require newer ones to be present first \citep{debianPolicyRelationships,debianPolicySharedLibraries}.  Although each record can look plausible in isolation, the risk appears only when the directed dependencies among records are followed.

\revredreplace{Large language models (LLMs) can read long contexts. Nominal context length alone does not guarantee reliable retrieval, multi-hop tracing, or aggregation under long inputs (\mbox{\citealp{liu2024lost}}; \mbox{\citealp{hsieh2024ruler}}; \mbox{\citealp{bai2024longbench}}).  Legal NLP benchmarks expose persistent difficulty with distributed evidence (\mbox{\citealp{hendrycks2021cuad}}; \mbox{\citealp{koreeda2021contractnli}}; \mbox{\citealp{chalkidis2022lexglue}}; \mbox{\citealp{guha2023legalbench}}).  We study a sharper failure mode in \emph{cyclic} document regions, where evidence is revisitable and diffuse.  A full-\scc{} one-shot prompt may find a suspicious node, but it is asked to simultaneously rank all records, produce a subgraph, and preserve multiple evidence endpoints in one generation.}{Large language models (LLMs) can read these records in a single context, but context length alone does not ensure reliable retrieval, multi-hop tracing, or evidence aggregation \citep{liu2024lost,hsieh2024ruler,bai2024longbench}. Following dependencies within strongly connected components (\scc{}s) can revisit the same records. One-shot prompting must jointly rank risk sources and assemble supporting evidence from the entire component in one generation. Complete evidence retention is therefore a distinct objective from locating a suspicious record.}

A natural alternative is AlphaGo-style search
\citep{silver2016mastering,kocsis2006bandit,browne2012survey}.  Standard MCTS unrolls trajectories as trees.  In a loopy graph, this makes the search state grow with
walk prefixes instead of physical records.  In an \scc{}, even a
small set of records can therefore induce many repeated path copies. Related MCTS analyses describe this as a transposition, graph-history, and loop-induced uncertainty
problem \citep{childs2008transpositions,kishimoto2004general,saffidine2012ucd,moerland2020mcts}.
\revredreplace{Prior Monte Carlo graph search reduces duplicate transpositions for
AlphaZero-style search by sharing states in an acyclic search graph
(\mbox{\citealp{czech2021monte}}).  Other graph-to-directed-acyclic-graph (DAG) strategies, such as removing
or orienting feedback edges to obtain an acyclic approximation
(\mbox{\citealp{eades1993fast}}; \mbox{\citealp{berger1990approximation}}), avoid cycles by discarding
part of the cyclic structure.  Neither line provides a mechanism for
turning a strongly connected document region into a stable evaluation unit.
Our setting requires localizing the \scc{} and collapsing it into an
evidence-preserving risk subgraph. The subgraph should retain explicit risk
sources and supporting evidence.}{Prior Monte Carlo graph search shares transpositions in acyclic search graphs \citep{czech2021monte}, while feedback-edge methods obtain directed acyclic graph (DAG) approximations by removing or reorienting edges \citep{eades1993fast,berger1990approximation}. Risk extraction requires a search representation that preserves cyclic relations while consolidating evidence about the same physical records.}

\revredreplace{We propose \textbf{\samcgs{}}, a structure-aware Monte Carlo graph search method for extracting risk subgraphs from cyclic document regions.  \samcgs{} first finds \scc{} bottlenecks, then performs local window rollouts inside each \scc{}.  Instead of treating risk as intrinsic to a single node, it stores relation-first evidence, where a node becomes risky when earlier evidence makes its constraint, exception, or handoff incompatible with the graph context.  The final output is a dynamic core risk subgraph that can be inspected or passed upward as a collapsed evaluation point.}{We propose \textbf{\samcgs{}}, a structure-aware Monte Carlo graph search method that performs local LLM rollouts within each \scc{}. A record\textquotesingle{}s priority depends on how its constraints, exceptions, or handoffs interact with accumulated evidence. Relation-first memory preserves this context across rollouts and guides exploration of implicated records and relations. Dynamic core replacement revises earlier selections as evidence accumulates, producing a compact risk subgraph for inspection (Figure~\ref{fig:pipeline}).}

Our contributions are:
\begin{itemize}[leftmargin=*,itemsep=2pt]
    \item \revredreplace{We formulate \emph{structural risk subgraph extraction}: given an \scc{} in a document graph, output a compact subgraph that retains the repair-relevant evidence endpoints, in place of single anomalous node.}{A formulation of \emph{structural risk subgraph extraction}: retaining repair-relevant evidence in compact subgraphs of cyclic document regions.}
    \item \revredreplace{We introduce \samcgs{}, a structure-aware Monte Carlo graph search method that localizes \scc{}s, accumulates relation-level evidence through local rollouts, and distills persistent signals into a compact dynamic risk core.  Its rollout policy uses upper confidence bound (UCB) selection augmented with progressive relation priors, critical-pair revisiting, and evidence-gated core pruning, so additional search budget is routed toward repeatedly repair-relevant relations, not arbitrary path copies.}{\samcgs{}, an adaptive search method combining upper confidence bound (UCB) selection with progressive relation priors, critical-pair revisiting, and evidence-gated core pruning.}
    \item \revredreplace{We build a frozen critical-stress benchmark from four authority-backed sources. Debian provides package-dependency graphs, SEC EX-21 provides subsidiary-disclosure graphs, the German Civil Code (BGB) provides statutory cross-reference graphs, and CUAD provides contract-clause graphs. The benchmark contains 320 model-case pairs evaluated with GPT-4o, DeepSeek-V3, Qwen2.5-72B, and Gemini 2.5 Pro. Across these settings, \samcgs{} improves \rootthree{}, \riskany{}, and \riskall{} over an oracle-risk full-\scc{} direct-subgraph baseline. It also makes the retention cost visible through an explicit compression metric.}{A fixed four-domain benchmark comparing \samcgs{} with one-shot prompting and LEA across five LLMs, with component ablations, inference-cost measurements, and expert evaluation. \samcgs{} achieves higher mean \rootthree{} and \riskall{} than either baseline.}
\end{itemize}

\begin{figure*}[!t]
    \centering
    \includegraphics[width=\textwidth]{fig01_sa_mcgs_framework_architecture_figma.png}
    \caption{\samcgs{} pipeline for \scc{} localization, local rollouts, relation-first evidence memory, and dynamic-core collapse.}
    \label{fig:pipeline}
\end{figure*}

\section{Task Formulation}
\label{sec:formulation}

\paragraph{Record graph.}
Let $G=(V,E)$ be a directed graph whose nodes are document records and whose edges denote typed dependencies such as reference, constraint, modification, trigger, or control.  A record may be a clause, statute section, package metadata entry, subsidiary disclosure, or other domain unit.  We use Tarjan's algorithm \citep{tarjan1972depth} to decompose $G$ into strongly connected components.  A component $S\subseteq V$ is a reasoning bottleneck when every node in $S$ can reach every other node through directed dependencies.

\paragraph{Structural risk.}
We study relational risks whose evidence extends beyond suspicious wording.  \revdel{In the benchmark, each injected critical case modifies only existing records and introduces no new nodes.}  The core evidence contains a root record, a witness record, and sometimes affected records.  The root and witness can each be locally plausible, but jointly create an incompatible obligation, temporal condition, ownership statement, or dependency handoff.

\paragraph{Output.}
Given a strongly connected component $S$, the system returns a risk subgraph $H\subseteq S$ and, when applicable, a ranking over risky records.  We evaluate:
\begin{itemize}[leftmargin=*,itemsep=2pt]
    \item \textbf{\rootthree{}}: whether the injected root appears in the top three ranked records.
    \item \textbf{\riskany{}}: whether $H$ contains at least one repair-relevant risk endpoint.
    \item \textbf{\riskall{}}: whether $H$ retains all core risk endpoints.  This criterion is strict since remediation often requires seeing both sides of the contradiction.
    \item \textbf{Compression}: $1-|H|/|S|$, where higher values indicate smaller subgraphs.  Compression is not optimized alone; a tiny subgraph that drops the evidence endpoints is not useful.
\end{itemize}
We report effective compression and effective \oc{} as diagnostics in Appendix~\ref{app:diagnostics}; they are not primary success metrics.
\revred{An unmodified component is a source \scc{}; each controlled edit produces an injected variant. Evaluating one variant with one model constitutes a model-case. Variants from the same source are not independent samples. Expert evidence and repair-material coverage are evaluated separately from the benchmark endpoint metrics.}

\section{Method}
\label{sec:method}

\subsection{Why ordinary MCTS is not enough}
\label{sec:mcts_motivation}

Standard tree search treats a trajectory as a growing path, following the usual MCTS loop of selection, expansion, simulation, and backpropagation \citep{browne2012survey}.  In a DAG, this assumption is valid. Rollouts eventually terminate, and backpropagation updates a finite set of states.  In an \scc{}, the same physical record can be revisited through many path copies \citep{childs2008transpositions}.  If a rollout path is $p=(v_1,\ldots,v_\ell)$, ordinary TreeMCTS indexes states by path prefix.  A physical node $v$ can therefore appear as many distinct tree states,
\begin{equation}
    \mathcal{P}_T(v)=
    \{p_{\leq t}: p_t=v,\; p \in \Pi_T\},
\end{equation}
where $\Pi_T$ is the set of paths generated before budget $T$.  With this indexing, path multiplicity governs the effective state space in place of $|S|$.  Repeated revisits then consume rollout budget, increase truncation, and dilute value estimates.

\revredreplace{Figure~\ref{fig:mcts_failure} illustrates the failure mode.  The issue is not simply that a cycle exists. Tree expansion also treats revisits as new path states.  A long cyclic document region can therefore consume rollout budget while repeatedly rediscovering the same records, making it harder to preserve the endpoints that jointly create a structural risk.}{Figure~\ref{fig:mcts_failure} contrasts path-indexed TreeMCTS with \scc{}-localized search. With matched rollout budgets, TreeMCTS expands 6.4 times as many states as physical nodes in the cyclic setting, and all rollouts are truncated. \samcgs{} limits expansion to 1.0 times the component size and increases Q-spread from 0.044 to 0.371. Table~\ref{tab:mcts_scc} includes the acyclic control; Appendix~\ref{app:diagnostics} defines the diagnostic setup and metrics.}

\begin{figure*}[!t]
    \centering
    \includegraphics[width=\textwidth]{fig02_vanilla_mcts_scc_failure_figma.png}
    \caption{Why vanilla TreeMCTS is ill-suited to \scc{}s. Cyclic records create path-copy states that dilute the rollout budget and flatten value estimates.}
    \label{fig:mcts_failure}
\end{figure*}

\revdel{We verify this failure mode with a controlled motivating ablation.  We take matched document-graph components with the same local risk task and compare three search variants under the same rollout budget: ordinary TreeMCTS on an acyclic dependency graph, ordinary TreeMCTS on an \scc{}, and \samcgs{} on the same \scc{} setting.  The ablation is intentionally small and diagnostic. It tests whether \scc{} topology breaks the assumptions behind tree search. Those assumptions underlie the appeal of tree search in DAG-like document graphs.}

% M1: retain the original diagnostic table in the main text, as requested.
\begin{table}[!h]
\centering
\scriptsize
\resizebox{\columnwidth}{!}{%
\begin{tabular}{llrrrr}
\toprule
Method & Topology & Expansion & Trunc. & Q-spread & Hit@3 \\
\midrule
TreeMCTS & DAG & 2.3$\times$ & 0\% & 0.077 & 7\% \\
TreeMCTS & SCC & 6.4$\times$ & 100\% & 0.044 & 30\% \\
\samcgs{} & SCC & 1.0$\times$ & -- & 0.371 & 90\% \\
\bottomrule
\end{tabular}
}
\caption{Motivating ablation: ordinary TreeMCTS expands path copies and dilutes value estimates inside strongly connected components, whereas \samcgs{} treats each \scc{} as a localized reasoning unit.}
\label{tab:mcts_scc}
\end{table}

\revdel{Table~\ref{tab:mcts_scc} directly tests the predicted failure. Inside an \scc{}, TreeMCTS expands many path-copy states, which truncates rollouts and flattens value estimates.  \samcgs{} avoids path copies by indexing evidence by physical nodes and relations, giving $1.0\times$ expansion and a much larger Q-spread.  Detailed definitions of Expansion, Truncation, Q-spread, and Hit@3 are given in Appendix~\ref{app:diagnostics}.}

\subsection{Objective}

\samcgs{} avoids path-copy states by treating each \scc{} $S$ as a localized graph object.  The desired output is a subgraph $H\subseteq S$ that keeps enough evidence for remediation while remaining smaller than the original component.  Let $R^\star\subseteq S$ denote the latent repair-relevant endpoints in evaluation, e.g., root and witness records.  At inference $R^\star$ is unknown, so the algorithm estimates node and relation utility from rollout evidence.  The empirical objective is given by
\begin{equation}
    \max_{H\subseteq S} \; \widehat{R}(H) - \lambda \frac{|H|}{|S|},
    \label{eq:objective}
\end{equation}
where $\widehat{R}(H)$ estimates evidence retention and $\lambda$ controls compression pressure.  The objective balances risk coverage, measured by \riskany{} and \riskall{}, against compression, reflecting the need for a small $H$ to retain repair-relevant evidence.

\subsection{SCC-aware local rollouts}

We first decompose the record graph using Tarjan's linear-time algorithm \citep{tarjan1972depth}.  DAG regions are handled by the ordinary upstream search, whereas each non-trivial \scc{} is processed by \samcgs{}.  At rollout $t$, the algorithm selects a seed $a_t\in S$ and constructs a local window
\begin{equation}
    W_t = \{v\in S : \operatorname{dist}_{S}(a_t,v)\leq r\} \cup B_t,
\end{equation}
where $r$ is a small graph radius and $B_t$ contains selected bridge or memory nodes.  The LLM evaluator observes node text and visible directed edges restricted to $W_t$, excluding the remainder of the \scc{}.  It returns local node scores $y_t(v)\in[0,1]$, relation evidence $z_t(u,v)\in[0,1]$, a local risk subgraph $h_t$, and optional \oc{} evidence.

Seed selection and frontier expansion follow the UCB intuition used in MCTS \citep{kocsis2006bandit}. Statistics are keyed by physical graph objects instead of path copies.  Let $x$ denote either a seed node $v$ or a frontier edge $e=(u,v)$.  \samcgs{} scores both with the same exploration template:
\begin{equation}
    \begin{aligned}
    U_t(x)=&\;\bar{r}_t(x)
    + c_t\sqrt{\frac{\log(1+N_t)}
    {n_t(x)+\ell_t(x)+\epsilon}}\\
    &+\frac{\lambda\pi_t(x)}
    {1+\delta\sqrt{n_t(x)}}+\eta_t(x),
    \end{aligned}
    \label{eq:ucb}
\end{equation}
where $\bar{r}_t(x)$ is the current empirical risk or conflict estimate, $n_t(x)$ is the number of times $x$ has been observed, $\ell_t(x)$ is virtual loss for concurrent workers, $\pi_t(x)$ is a relation-first prior, and $\eta_t(x)$ is exploration noise.  The prior is computed from rollout evidence without domain-specific labels.
\begin{equation}
    \pi_t(x)=
    \frac{[w^\top\phi_t(x)]_+}
    {\max_{x'}[w^\top\phi_t(x')]_+ + \epsilon}.
    \label{eq:relation-prior}
\end{equation}
For nodes, $\phi_t(x)$ contains repair-entry counts, local risk-subgraph appearances, conflict-pair endpoint support, AMAF/RAVE-style evidence, path-fragment support, and \oc{} endpoint indicators.  For edges, it contains semantic incompatibility, repeated pair support, explicit conflict reports, reverse-edge support, and accumulated conflict evidence.  Weights $w$ are fixed before the locked evaluation and reported in Appendix~\ref{app:priors}.  The decay term in Eq.~\ref{eq:ucb} prevents an early prior from dominating after sufficient node or edge samples.  Together, these choices implement progressive relation bias, analogous in spirit to progressive bias and RAVE-style reuse \citep{gelly2007combining,chaslot2008progressive,browne2012survey}, while keeping the update key tied to physical nodes and edges as graph-level objects \citep{czech2021monte}.

\subsection{Relation-first evidence memory}

Structural risk is relational in the sense that a node can appear harmless until previous evidence makes it incompatible with the graph context.  This matches the relational-inductive-bias view that graph-structured reasoning should represent entities together with the relations that define their roles \citep{battaglia2018relational}.  We therefore maintain memory
\begin{equation}
    M_t = \{n_t(v),\bar{q}_t(v),m_t(u,v),o_t(v)\}_{v,u\in S},
\end{equation}
where $m_t(u,v)$ stores support that the directed relation $(u,v)$ makes either endpoint risky, and $o_t(v)$ stores whether $v$ has appeared in effective \oc{} evidence.  Node evidence is updated by running averages and relation evidence by smoothed pair support. Appendix~\ref{sec:appendix_formula_details} gives the exact update equations.  Let $r^-_t(v)$ and $r^+_t(v)$ be the strongest incoming and outgoing relation-memory scores around $v$.  The priority score used for core construction combines node evidence, relation evidence, \oc{} support, and a persistent critical-pair ledger:
\begin{equation}
    \begin{aligned}
    s_t(v)=&\;\bar{q}_t(v)
    + \alpha r^-_t(v)+\alpha r^+_t(v)\\
    &\quad
    + \gamma o_t(v)+\mu p_t(v).
    \end{aligned}
    \label{eq:priority}
\end{equation}
\revdel{Thus a node can become important because incoming or outgoing evidence repeatedly implicates it, even if its isolated text is not obviously defective.}  Relation evidence is aggregated by maxima over incoming and outgoing edges to limit degree bias from highly connected nodes.  \revredreplace{The bounded bonuses from \oc{} and the pair ledger route search toward severe recurring hypotheses, allowing persistent node or relation evidence to outweigh isolated local judgments.  This persistence mechanism prevents a single noisy local judgment from deciding the final core, while still letting strong recurring relations guide pruning.}{Bounded \oc{} and pair-ledger bonuses prioritize severe recurring hypotheses, allowing persistent relation evidence to outweigh isolated local judgments.}

\subsection{Online conformal-style evidence signal}

Single LLM judgments are noisy, especially on local windows.  We therefore add a calibration-inspired discovery signal.  At each rollout, \samcgs{} compares the strongest local relation-conflict score against an online quantile of previous conflict scores. Appendix~\ref{sec:appendix_formula_details} gives the exact rule.  We call a threshold crossing an online calibrated conflict signal (\oc{}). The signal marks a local window whose ordered cycle or contradiction cue is unusually strong relative to earlier rollouts, and endpoints of the implicated relation receive an \oc{} bonus in Eq.~\ref{eq:priority}.  The calibration motivation follows conformal calibration \citep{vovk2005algorithmic,gibbs2021adaptive}. No formal coverage guarantee is claimed because the LLM evaluator and rollout policy are adaptive.  We use \oc{} as evidence routing and convergence support, while final evaluation is still based on the returned subgraph $H$.

\subsection{Dynamic core compression}
\label{sec:dynamic-core}

We maintain the final risk subgraph as a dynamic core with replacement, avoiding monotonic expansion.  Core compression is cast as budgeted subset selection, for which greedy selection is a standard approximation strategy under monotone submodular objectives with cardinality constraints \citep{nemhauser1978analysis}.  After each rollout, \samcgs{} forms a candidate pool $C_t$ from the local risk subgraph, the previous core, \oc{} endpoints, and top-scoring relation endpoints with local neighbors.  The core is selected by maximizing a bounded utility:
\begin{equation}
    \begin{aligned}
    H_t &= \arg\max_{\substack{H\subseteq C_t\\ |H|\leq B(S)}}
    \Bigg[
    \sum_{v\in H}s_t(v)\\
    &\quad + \eta \sum_{(u,v)\in E(H)}m_t(u,v)\\
    &\quad - \lambda |H|
    \Bigg].
    \end{aligned}
    \label{eq:core}
\end{equation}
We implement this objective with greedy replacement: high-confidence endpoints may evict weaker nodes, whereas repeatedly supported relation endpoints are retained as locked evidence.  \revreplace{This persistence mechanism is essential for critical cases.}{This design aims to retain critical evidence under a fixed core budget.} When a rollout observes both endpoints of a severe contradiction, subsequent rollouts should revisit and stabilize that relation, instead of allowing unrelated high-scoring distractors to occupy the subgraph. \revred{The paired replay in Section~\ref{sec:component_tests} isolates the effect of final-core replacement using identical search trajectories.}

\begin{algorithm}[t]
\small
\caption{\samcgs{} for one strongly connected component}
\label{alg:samcgs}
\begin{algorithmic}[1]
\REQUIRE Component $S=(V_S,E_S)$, budget $T$, shared rubric and domain serializer
\STATE Initialize memory $M_0$ and core $H_0=\emptyset$
\FOR{$t=1$ to $T$}
    \STATE Select seed $a_t=\arg\max_{v\in V_S}U_t(v)$ using Eq.~\ref{eq:ucb}
    \STATE Build local window $W_t$ around $a_t$ plus memory bridges
    \STATE Query LLM with the shared risk rubric on $W_t$
    \STATE Receive $y_t(v)$, $z_t(u,v)$, local subgraph $h_t$, and \oc{} evidence
    \STATE Update node and relation memory $M_t$
    \STATE Recompute dynamic core $H_t$ using Eq.~\ref{eq:core}
\ENDFOR
\RETURN risk subgraph $H_T$ and node ranking by $s_T(v)$
\end{algorithmic}
\end{algorithm}

\subsection{Computational feasibility}

\samcgs{} is an anytime graph-search procedure in the sense of
\citet{zilberstein1996anytime}.  Tarjan preprocessing is
$O(|V|+|E|)$.  For a component $S$, the method performs $T$ local LLM
evaluations over windows of at most $w$ records, with prompt cost
$O(T\cdot\operatorname{tokens}(w))$.  Its evidence memory is keyed by
physical nodes and observed relations, not path-copy states, and stores
$O(|S|+|E_S|)$ statistics.  The dynamic core cap $B(S)$ bounds the
retained subgraph, allowing additional budget to refine the same physical
memory without expanding a cyclic tree indefinitely.  The resulting
collapsed evaluation point is anytime-compatible for higher-level graph search.
\revredreplace{Budget-prefix behavior is analyzed in Figure~\ref{fig:convergence} and Appendix~\ref{app:diagnostics}.}{Appendix~\ref{app:budget_prefix} analyzes performance across rollout-budget prefixes.}
\revred{Section~\ref{sec:cost_analysis} reports call and token counts; concurrent rollouts can reduce waiting time without reducing the inference workload.}

\section{Experimental Setup}
\label{sec:setup}

\subsection{Domains}

\revnewreplace{We evaluate on four domains chosen for structural diversity and source authority.  The selection was not designed to favor a particular result.  Table~\ref{tab:domains} summarizes their roles.}{We evaluate on Debian dependencies, SEC EX-21 subsidiary disclosures, German Civil Code (BGB) cross-references, and CUAD contracts, chosen for structural diversity and source authority. Appendix~\ref{sec:appendix_datasets} summarizes their roles and graph construction.}

% M2: retain the original black table only as deletion markup; merge its
% coverage information into the dataset-processing table in Appendix I.
\revdeletedtable{
\begin{table*}[t]
\centering
\small
\begin{tabular}{p{0.13\textwidth}p{0.24\textwidth}p{0.27\textwidth}p{0.25\textwidth}}
\toprule
\revdel{Domain} & \revdel{Source authority} & \revdel{Graph role} & \revdel{Paper role} \\
\midrule
\revdel{Debian} & \revdel{Package dependency metadata} & \revdel{Software migration / dependency constraints} & \revdel{Non-legal technical generalization} \\
\revdel{SEC EX-21} & \revdel{Public subsidiary exhibit structure} & \revdel{Entity-control and consolidation paths} & \revdel{Financial/legal entity generalization} \\
\revdel{BGB} & \revdel{German Civil Code references} & \revdel{Statutory cross-reference cycles} & \revdel{Clean legal-rule stress test} \\
\revdel{CUAD} & \revdel{Expert contract review dataset} & \revdel{Long contract clause graph with ambiguity} & \revdel{Noisy contract-domain stress test} \\
\bottomrule
\end{tabular}
\caption{\revdel{Domain coverage.  The benchmark spans technical dependencies, public regulatory disclosures, official statutory references, and expert-labeled commercial contracts.}}
\label{tab:domains:previous}
\end{table*}
}

\subsection{Domain formatting and shared risk semantics}

\ifrevisionmarks
\revdel{The four sources differ in surface form. A Debian record exposes package fields, a BGB record is a statutory section, an SEC EX-21 record is a subsidiary disclosure, and a CUAD record is a contract clause.  We therefore use lightweight domain formatters}
\begin{equation*}
\text{\revdel{\mbox{$\mathcal{A}_d=(\operatorname{serialize}_d,\operatorname{edges}_d,\operatorname{normalize}_d)$}}}
\end{equation*}
\revdel{to serialize records, expose typed directed edges, and normalize model outputs.  These formatters do not define the target risk type and do not reveal the injected endpoints.  Both \naive{} and \samcgs{} use the same risk rubric that is neutral across domains.  Under this rubric, models flag structural inconsistency or mutual incompatibility, as well as risks that are underspecified or have high impact, using only visible node text and directed edges.  The comparison is therefore computational, contrasting full-\scc{} direct-subgraph generation and Top-3 aggregation with \scc{}-aware rollout selection, relation-first memory, critical-pair revisiting, and dynamic-core compression.}
\fi

\revred{All methods receive the same visible record text and directed edges under a shared risk rubric covering structural inconsistency, mutual incompatibility, underspecification, and high-impact risks. Domain formatters preserve these risk semantics without revealing injected endpoints; Appendix~\ref{sec:appendix_datasets} details their interface and graph construction.}

\subsection{Critical structural injection}

\revredreplace{The main experiment uses a locked critical-stress configuration over the same \scc{} blocks and models for both methods. The exact run profile is reported in Appendix~\ref{sec:appendix_experiments}.  Each case edits only existing records in a real \scc{}.  A root record and a distant witness or affected record jointly create a structural conflict, while each record remains locally plausible.  Both methods receive the same generic risk semantics and the same visible graph; the exact templates and domain examples are in Appendix~\ref{sec:appendix_datasets}.}{Each injected variant modifies existing records in a real \scc{} without adding nodes. The 20 source \scc{}s each contribute four variants, yielding the same 80 cases per model. Five models and three methods produce 400 model-cases and 1,200 outputs. Inputs, injection identities, prompt schemas, and method settings are fixed throughout. Appendix~\ref{sec:appendix_datasets} provides the injection templates, and Appendix~\ref{app:revision_lea} specifies the run configuration.}

\subsection{Baselines and models}

\revdel{We use oracle-risk \naive{} as the main baseline. It runs the full-\scc{} direct-subgraph prompt ten times. From these attempts, a risk-prioritized oracle chooses the reported Top-3 set using \rootthree{}, \riskany{}, and \riskall{}. Each \naive{} attempt sees the full \scc{} and emits a whole-\scc{} ranking plus a risk subgraph; Appendix~\ref{sec:appendix_experiments} reports non-oracle and compression-first variants. \samcgs{} instead sees smaller local windows over 60 rollouts and aggregates evidence into a dynamic core.}
\input{revision_baselines}

\revredreplace{We evaluate GPT-4o, DeepSeek-V3, Qwen2.5-72B-Instruct, and Gemini 2.5 Pro.}{We evaluate \samcgs{} with five LLM backbones: proprietary GPT-4o~\citep{openai2024gpt4o} and Gemini 2.5 Pro~\citep{comanici2025gemini25}, and open-weight Qwen2.5-72B~\citep{qwen2024qwen25}, Llama 3.1-8B-Instruct~\citep{grattafiori2024llama3}, and Qwen3.5-9B~\citep{qwen2026qwen35nine}.}  \revdel{The main comparison contains 320 model-case pairs per method; oracle-risk \naive{} is selected from 3,200 raw full-\scc{} attempts.  This gives boosted \naive{} a total token budget comparable to \samcgs{}.}

\section{Main Results}
\label{sec:results}

\input{revision_main_results}



\revdel{Figure~\ref{fig:main_metrics} shows the core result.  Against oracle-risk \naive{}, \samcgs{} improves \rootthree{} by 23 points, \riskany{} by 19 points, and \riskall{} by 16 points.  The strongest claim is not that full-\scc{} prompting never finds risk, since the baseline is already selected from ten attempts with access to risk metrics.  Improvement from \samcgs{} is more reliable preservation of the repair-relevant evidence endpoints in the final subgraph.}

\revredreplace{The compression axis should be read as a trade-off, not a loss.  Oracle-risk \naive{} is slightly more compressed on average, but still retains all critical endpoints less often.  \samcgs{} deliberately keeps a larger dynamic core when evidence persists across rollouts.}{Paired \riskall{} outcomes favor \samcgs{} over LEA for GPT-4o, Qwen2.5-72B, and Llama 3.1-8B-Instruct (Figure~\ref{fig:paired_case_outcomes}b). For GPT-4o, 32 cases succeed only with \samcgs{}, compared with 7 only with LEA.}

\revdel{Appendix~\ref{tab:reliability} separates one-shot output availability from boosted raw-attempt and model-case reliability.}


\begin{revnewblock}
\begin{figure*}[t]
    \centering
    \includegraphics[width=\textwidth]{fig03_paired_case_outcomes_five_models.pdf}
    \caption{\revred{Paired \riskall{} outcomes for \samcgs{} versus (a) One-shot and (b) LEA on the same 80 cases per model. Bars separate cases where complete evidence is retained by \samcgs{} only (blue), both methods (light blue), the baseline only (orange), or neither method (gray). Net $\Delta$ is the \riskall{} difference between \samcgs{} and the baseline, in percentage points.}}
    \label{fig:paired_case_outcomes}
\end{figure*}
\end{revnewblock}

\subsection{Breakdown by model and domain}

\revredreplace{The gains are not carried by a single model.  Gemini 2.5 Pro and DeepSeek-V3 are strong even under oracle-risk \naive{}, while GPT-4o and Qwen2.5-72B still lag in complete endpoint retention. Domain-wise, CUAD is the hardest setting: long contract components create both context pressure and output-schema burden for full-\scc{} ranking and subgraph generation.}{The domain breakdown also shows gains in complete evidence retention (Table~\ref{tab:appendix_breakdowns}, panel B). Pooling results across the five models, \samcgs{} raises \riskall{} over LEA from 26.67\% to 51.67\% on BGB and from 53.00\% to 69.00\% on CUAD. These gains extend to statutory cross-references and longer contract-clause graphs.}
\revred{Appendix~\ref{sec:appendix_case} compares the three methods on the same long-\scc{} case.}

\revred{Across all three \scc{}-size groups, \samcgs{} improves \riskall{} over LEA for GPT-4o and Qwen2.5-72B, and improves \rootthree{} for Llama 3.1-8B-Instruct (Appendix~\ref{app:size_breakdowns}).}

\section{Analysis}
\label{sec:analysis}

% October 1: full size analysis moved to Appendix J.1; original black text retained.
\ifrevisionmarks
\subsection{\revdel{Effect of SCC size}}
\fi
\revdel{Figure~\ref{fig:main_metrics}B separates the results by \scc{} size.  In short components ($\leq 12$ nodes), \samcgs{} improves \riskall{} from 79\% to 83\%.  In mid-sized components (14--20 nodes), the improvement is 47\% to 80\%.  In long components ($\geq 24$ nodes), \riskall{} remains difficult, but \samcgs{} still improves 50\% to 65\%.  The pattern is consistent with the task. In longer cycles, plausible distractors are more numerous and native ambiguity is higher, especially in CUAD.}

\subsection{Rollout convergence}

\revredreplace{Figure~\ref{fig:convergence} shows that additional rollouts accumulate useful evidence instead of only spending more tokens.  At budget 5, \riskall{} is 22\%; at budget 30, it rises to 41\%; at budget 60, it reaches 57\% under prefix reconstruction.  The final run-level \riskall{} in Figure~\ref{fig:main_metrics} is higher because the dynamic core uses the full retained evidence rather than a simple prefix truncation.  The important trend is that risk coverage and \riskall{} continue to rise well past the first few rollouts, supporting the search interpretation.}{For GPT-4o, most gains in evidence completeness occur within the first 20 rollouts; subsequent rollouts refine the retained core (Appendix~\ref{app:budget_prefix}).}

\ifrevisionmarks
\subsection{\revdel{Robustness checks}}
\fi

\revdel{Appendix~\ref{sec:appendix_experiments} reports selector and reliability checks.  The non-oracle boosted selector reaches 52\% \riskall{}, while compression-first oracle selection raises compression but lowers \riskall{} to 47\%.  A CUAD lite-output control reduces formatting burden but still reaches only 12\% \riskall{} on long-contract SCCs, and a more compressed \samcgs{} profile lowers \riskall{} from 78\% to 71\%.
Appendix~\ref{sec:appendix_case} gives a representative long-\scc{} case study.}

\input{revision_main_analysis}

\section{Related Work}
\label{sec:related}

\paragraph{LLM reasoning and long context.}
Chain-of-thought prompting improves sequential reasoning \citep{wei2022chain}, whereas cyclic document risk requires revisiting and reconciling records instead of extending one chain.  Long-context benchmarks show persistent failures in retrieval, multi-hop tracing, and aggregation \citep{liu2024lost,hsieh2024ruler,bai2024longbench}, while graph-reasoning studies show sensitivity to graph encoding and task formulation \citep{wang2023can,fatemi2024talk}.  Legal NLP benchmarks motivate document-level reasoning \citep{hendrycks2021cuad,koreeda2021contractnli,chalkidis2022lexglue,guha2023legalbench}. Our task stresses cyclic structural evidence retention.

\paragraph{Search-augmented LLMs.}
Tree of Thoughts \citep{yao2024tree}, Graph of Thoughts \citep{besta2024graph}, and planning-style LLM search \citep{hao2023reasoning} show that explicit search can improve reasoning.  \samcgs{} operates on externally constructed document graphs and treats \scc{}s as bottlenecks to collapse into risk subgraphs.

\paragraph{Monte Carlo search on graphs.}
MCTS and AlphaGo-style search are powerful under stable state transitions \citep{silver2016mastering,kocsis2006bandit,browne2012survey}.  Monte Carlo graph search reuses transpositions \citep{childs2008transpositions,czech2021monte}, and work on MCTS uncertainty highlights loop difficulty \citep{moerland2020mcts}.  Our method targets the remaining gap in information-preserving \scc{} handling for cyclic document regions.

\paragraph{Graph legal analysis and calibration.}
Network analysis has been used for statutory cross-references \citep{boulet2021measuring}, contract graph extraction is active \citep{dechtiar2025graphgrpolex}, and conformal ideas inspire calibrated evidence signals \citep{vovk2005algorithmic,gibbs2021adaptive}.  \samcgs{} is complementary because it assumes a record graph and focuses on risk search inside cyclic components.

\section{Conclusion}
\label{sec:conclusion}

We introduced \samcgs{}, a structure-aware Monte Carlo graph search method for cyclic document risk.  \revredreplace{The central idea is to avoid treating a long \scc{} as either a flat prompt or an ordinary tree-search path.  Instead, \samcgs{} localizes the \scc{}, searches it through local rollouts, stores relation-first evidence, and collapses the component into a dynamic risk core.  Across four domains and four LLMs, this yields stronger root localization and endpoint retention than an oracle-risk Top-3 baseline built from repeated full-\scc{} direct-subgraph attempts. It also makes the cost in compression explicit.}{Across four domains and five LLMs, \samcgs{} achieves higher mean \rootthree{} and \riskall{} than one-shot prompting and LEA. Its compact risk subgraphs consolidate cyclic dependencies into evidence for human inspection.}

\section*{Limitations}
\label{sec:limitations}

\samcgs{} incurs higher cost than one-shot prompting due to its local LLM calls.  The method is intended for high-risk cyclic components, not every record in a document.  Our benchmark uses injected critical risks to keep ground truth controlled.  \revdel{Naturally occurring risks may be less controlled and require human audit.}\revred{The 18 audited unmodified sources are expert-confirmed negatives, allowing us to assess over-reporting but not recall on naturally occurring risks or domain-wide prevalence. Stage-1 covers a small SEC subset, and the sensitivity analysis varies one prior coefficient on four sources. These designs do not establish uniform component necessity or superiority at equal token budgets.} Graph construction quality also matters, as missing or spurious edges can change the \scc{}s that the method receives.  Finally, \riskall{} remains difficult in noisy long-contract settings, so we report it explicitly without replacing it with the easier \riskany{} metric. \revred{False positives and missing evidence require human review; downstream decisions over complete documents remain unevaluated.}

\bibliography{custom}

\appendix
\raggedbottom


\section{Diagnostic Metrics and Ablation Definitions}
\label{app:diagnostics}

The main paper evaluates \rootthree{}, \riskany{}, \riskall{}, and
compression.  We also log two diagnostics.  \textbf{Effective
compression} is computed only on cases whose returned subgraph retains at
least one repair-relevant endpoint.  It separates useful compression from
trivial compression that drops all evidence.  \textbf{Effective \oc{}}
counts an \oc{} discovery only when the implicated relation intersects
the repair-relevant risk region.  It is useful for convergence analysis,
although it is cumulative by construction and is therefore not a primary
success metric.

\begin{revredblock}
\paragraph{Motivating search diagnostic.}
We compare TreeMCTS on acyclic and strongly connected document graphs with \samcgs{} on the same \scc{} setting. The comparison uses matched components, the same local risk task, and equal rollout budgets. Table~\ref{tab:mcts_scc} reports state expansion, rollout truncation, Q-spread, and Hit@3, as defined below.


\end{revredblock}

Table~\ref{tab:mcts_scc} uses four ablation diagnostics.  \textbf{Expansion}
is the number of search states expanded relative to the physical
component size.  \textbf{Truncation} is the fraction of rollouts stopped
by the depth limit before a terminal state, which is common in loops
\citep{moerland2020mcts}.  \textbf{Q-spread} is the range of estimated
action values at the root, a proxy for whether search can separate
promising actions under UCT-style selection \citep{kocsis2006bandit}.
\textbf{Hit@3} reports whether a repair-relevant endpoint appears in the
top three ranked records.

\section{System-Blind Expert Annotation Audit}
\label{sec:appendix_annotation}

\revred{We report the original 12-case audit below and the new 42-presentation audit in Section~\ref{app:expert_extension} separately because their annotation schemes and annotator coverage differ. The original aggregate results are retained without imputing missing labels.}
\subsection{\revadd{Original 12-case audit}}

We construct a representative expert-audit set to verify that the critical structural risks used in our benchmark are recognizable to qualified domain readers, not artifacts of hidden system labels.  The audit is \emph{system-blind}.  Annotators see only paragraph text, paragraph identifiers, and ordinary natural-language instructions.  They are not shown model outputs, hidden node roles, system suggestions, \samcgs{} cores, \revnewreplace{\naive{} subgraphs}{baseline subgraphs}, or terms such as root, witness, affected node, \scc{}, and risk subgraph.  Hidden fields are used only after annotation to compare expert labels with injected endpoints and system outputs.

All three annotators hold doctoral degrees in relevant fields and have professional experience with legal, regulatory, or contractual analysis. \revred{Table~\ref{tab:expert_protocol} summarizes the annotation tasks.}

\paragraph{Case-level task.}
For each case, the expert reads a group of related paragraphs and answers whether the group contains a structural conflict.  A structural conflict means that two or more paragraphs cannot all be satisfied together, leave a high-impact condition unresolved, or make a repair decision depend on inconsistent assumptions.

\paragraph{Paragraph-level task.}
For each paragraph, the expert chooses one of four labels:
\begin{itemize}[leftmargin=*,itemsep=1pt]
    \item \textbf{Clear issue}: the paragraph is itself problematic, or clearly conflicts with another visible paragraph.
    \item \textbf{Possibly relevant}: the paragraph may be needed to understand or repair the conflict, but is not clearly wrong by itself.
    \item \textbf{No issue observed}: the paragraph appears to be ordinary context.
    \item \textbf{Uncertain}: the expert cannot decide without more context or specialist knowledge.
\end{itemize}

\paragraph{Repair-material task.}
After labeling paragraphs, the expert selects which paragraphs should be included in the final repair packet.  This task is stricter than identifying a suspicious sentence because the packet should include enough material for a reviewer to understand and fix the conflict.

\paragraph{Annotation protocol.}
Each sample is annotated independently by three experts.  Every sample therefore receives three case-level decisions, three sets of paragraph-level labels, and three repair-packet selections.  Disagreements are preserved for analysis without being forced into an immediate majority vote.

\paragraph{Audit results.}
Table~\ref{tab:expert_audit_results} summarizes the completed annotations.  All 12 audited cases receive a majority structural-conflict label.  Experts also recognize the injected evidence endpoints.  In total, 23 of 24 root/witness endpoint paragraphs receive a majority \emph{clear issue} label, and all 24 receive a majority \emph{clear issue} or \emph{possibly relevant} label.  Agreement is strongest for the clear-risk endpoint decision, with Fleiss' $\kappa=0.606$ and pairwise Cohen's $\kappa$ ranging from 0.452 to 0.759.  These results support that the injected critical risks are semantically visible to qualified readers when the relevant paragraphs are inspected.

The repair-packet task is intentionally stricter and more variable.  One annotator selected no individual paragraph for the final packet while still marking every case as conflicting and most cases as sufficiently or partially supported.  We therefore treat repair-packet selection as a conservative diagnostic of what a reviewer would include for remediation, not as the primary evidence that a conflict exists.

\begin{table*}[t]
\centering
\small
\setlength{\tabcolsep}{4pt}
\begin{tabular}{p{.22\textwidth}p{.20\textwidth}p{.50\textwidth}}
\toprule
Audit quantity & Result & Interpretation \\
\midrule
Audited critical cases & 12 cases, 36 case-level decisions & Each case is independently labeled by three experts. \\
Majority structural conflict & 12/12 (100\%) & The audited critical cases are judged to contain structural conflicts. \\
Repair material sufficiency & 10 yes, 1 partial, 1 tie & Most cases contain enough visible material to support repair, but one case remains divided. \\
No extra context needed & 10/12 (83.3\%) & Most cases are understandable from the provided paragraphs alone. \\
Mean confidence & 4.28/5 & Experts are generally confident in their case-level judgments. \\
Endpoint majority clear issue & 23/24 (95.8\%) & Root/witness endpoints are usually judged directly problematic. \\
Endpoint majority clear-or-relevant & 24/24 (100\%) & Every endpoint is recognized as at least repair-relevant. \\
Clear-issue agreement & Fleiss' $\kappa=0.606$ & Agreement is substantial for the binary endpoint-risk judgment. \\
Repair-packet endpoint-any & 7/12 (58.3\%) & Majority repair packets include at least one endpoint in over half of cases. \\
Repair-packet endpoint-all & 6/12 (50.0\%) & Selecting all endpoints is stricter and less stable across experts. \\
\bottomrule
\end{tabular}
\caption{System-blind expert audit results.  Experts did not see system outputs, hidden roles, or model-generated hints during annotation.}
\label{tab:expert_audit_results}
\end{table*}

\paragraph{Role-level sanity check.}
After annotation, we compare expert paragraph labels with hidden roles.  True endpoints are the only role class with near-universal majority recognition, with 95.8\% majority clear issue and 100\% majority clear-or-relevant.  In contrast, affected paragraphs are rarely majority clear issue, and most broader system-core paragraphs are judged ordinary context.  This behavior follows from the benchmark design, which asks methods to retain compact endpoints needed for repair, not to mark every retained context paragraph as intrinsically wrong.  The same observation explains why \riskall{} is strict, as a method may surface useful context or affected nodes while still missing one side of the contradiction.

\begin{table*}[t]
\centering
\small
\setlength{\tabcolsep}{4pt}
\begin{tabular}{p{.15\textwidth}p{.33\textwidth}p{.42\textwidth}}
\toprule
Task & Expert-facing question & Example expert note \\
\midrule
Case-level conflict & Do these paragraphs create a conflict or unresolved high-risk condition when read together? & ``The first paragraph gives a right that appears unconditional, but a later paragraph makes the required approval unavailable.  The group needs revision before the right can be applied.'' \\
Paragraph label & Does this paragraph look clearly problematic, possibly relevant, ordinary, or uncertain? & ``This paragraph is not wrong alone, but it is relevant because it provides the timing condition that controls the later obligation.'' \\
Repair packet & Which paragraphs should a reviewer see to fix the issue? & ``The packet should include the grant clause, the later condition clause, and the definition paragraph.  Omitting the definition would make the repair ambiguous.'' \\
\bottomrule
\end{tabular}
\caption{Expert-facing annotation tasks.  These are the instructions used to elicit human judgments; internal roles and system outputs are used only for post-hoc analysis.}
\label{tab:expert_protocol}
\end{table*}

\begin{table*}[t]
\centering
\small
\setlength{\tabcolsep}{4pt}
\begin{tabular}{p{.16\textwidth}p{.50\textwidth}p{.24\textwidth}}
\toprule
Displayed paragraph & Expert-visible text excerpt & Majority annotation \\
\midrule
SEC EX-21 entity record A & \emph{Entergy Louisiana Holdings LLC.} ``For this consolidation path, the entity is treated as fully consolidated by the upstream registrant.  The note states that no separate non-controlling interest is retained for this path.'' & Clear issue; included in two experts' repair packet selections. \\
SEC EX-21 entity record B & \emph{ENTERGY LOUISIANA HOLDINGS INC.} ``The same consolidation path reserves a 20 percent non-controlling interest for the downstream schedule.  This record applies the opposite state immediately on the same path.'' & Clear issue; included in two experts' repair packet selections. \\
Case-level decision & The two records describe the same consolidation path and assign incompatible ownership states, fully consolidated with no non-controlling interest versus a retained 20\% non-controlling interest. & 3/3 experts mark the case as structurally conflicting; 2/3 mark the provided material sufficient, and 1/3 marks it partially sufficient. \\
\bottomrule
\end{tabular}
\caption{Representative expert-audit example. \revredreplace{Role names in the left column are added here for exposition; the annotators saw paragraph text and identifiers, not root/witness labels or system suggestions.}{Annotators saw paragraph text and identifiers, without hidden root/witness labels or system suggestions.}}
\label{tab:expert_example}
\end{table*}

The example in Table~\ref{tab:expert_example} illustrates the intended difficulty level.  Neither entity record is a syntactic error by itself.  The conflict appears because both records refer to the same consolidation path while asserting mutually incompatible ownership treatments.  A risk-subgraph method should preserve this type of structural conflict because a reviewer can reliably repair the issue after seeing both endpoints.

\input{revision_expert_appendix}

% Keep each appendix's floats before the next unrelated appendix.
\FloatBarrier
\section{Representative Long-\scc{} Case Study}
\label{sec:appendix_case}

\begin{figure*}[t]
    \centering
    \includegraphics[width=\textwidth]{\revasset{fig06_representative_scc_collapse.pdf}{fig06_representative_scc_three_methods.pdf}}
    \caption{\revredreplace{Representative long-\scc{} collapse.  The method searches a large component through local rollouts, accumulates relation-first evidence, and returns a compact risk subgraph that preserves repair-relevant endpoints.}{Comparison of three methods on the same 25-record BGB case with Gemini 2.5 Pro. Node positions are fixed across panels; filled markers indicate retained records and open markers indicate omitted records. Diamonds mark the root (\S505) and witness (\S491). The rows below report retained node counts and the four evaluation metrics.}}
    \label{fig:case}
\end{figure*}

\revredreplace{Figure~\ref{fig:case} illustrates the intended use of the method.  The original \scc{} is too large to treat as a single atomic finding and too recurrence-heavy for ordinary path search.  \samcgs{} localizes the evidence through rollouts, then returns a core containing the root, witness, and affected records.  For review, this is more useful than a single node label because a human can inspect the retained endpoints and decide how the conflict should be repaired.}{Figure~\ref{fig:case} compares evidence selection within the same cyclic region. One-shot retains three records but omits the root. LEA and \samcgs{} retain all risk endpoints, but \samcgs{} uses nine records compared with twelve for LEA, increasing compression from 52\% to 64\%. The shared layout allows direct comparison of the retained constraint and its distant witness.}

\FloatBarrier
\section{Core Algorithm Details}
\label{sec:appendix_algorithm}

\revredreplace{This appendix expands the method into reproducible pseudocode.}{Algorithms~\ref{alg:appendix_scc_localized}--\ref{alg:appendix_dynamic_core} expand the overview in Algorithm~\ref{alg:samcgs} into detailed pseudocode.}  The implementation uses the same high-level objects as the main paper, including a document graph $G$, strongly connected components $S$, local rollout windows $W_t$, relation-first memory $M_t$, online calibrated conflict signals, and a final dynamic core $H_T$.

\begin{algorithm}[H]
\small
\caption{\scc{}-localized search}
\label{alg:appendix_scc_localized}
\begin{algorithmic}[1]
\REQUIRE Document graph $G=(V,E)$, rollout budget $T$, shared risk rubric
\STATE $\mathcal{S}\leftarrow$ Tarjan strongly connected components of $G$
\STATE Initialize collapsed graph $G' \leftarrow G$
\FORALL{component $S\in\mathcal{S}$}
    \IF{$|S|=1$}
        \STATE Keep the singleton record unchanged
    \ELSE
        \STATE Build local search unit $G_S=(S,E_S)$
        \STATE $(H_S,r_S)\leftarrow \textsc{SA-MCGS-Rollout}(G_S,T)$
        \STATE Replace $S$ in $G'$ with collapsed evaluation point $(H_S,r_S)$
    \ENDIF
\ENDFOR
\RETURN Collapsed graph $G'$ and all component risk subgraphs
\end{algorithmic}
\end{algorithm}

\begin{algorithm}[H]
\small
\caption{\samcgs{} rollout inside one component}
\label{alg:appendix_rollout}
\begin{algorithmic}[1]
\REQUIRE Component $G_S=(S,E_S)$, budget $T$, shared rubric
\STATE Initialize node evidence, relation evidence, \oc{} evidence, and core $H_0=\emptyset$
\FOR{$t=1$ to $T$}
    \STATE Select a seed record using SCC-aware UCB with relation priors
    \STATE Construct local window $W_t$ from the seed, nearby edges, and memory bridges
    \STATE Query the LLM with the shared risk rubric on $W_t$
    \STATE Parse local risky records, relation evidence, and optional \oc{} signal
    \STATE Update node risk scores, pair evidence, witness-style evidence, and affected-region evidence
    \STATE $H_t\leftarrow \textsc{Dynamic-Core-Compression}(H_{t-1},M_t)$
\ENDFOR
\RETURN Final core $H_T$ and ranking induced by accumulated evidence
\end{algorithmic}
\end{algorithm}

\begin{algorithm}[H]
\small
\caption{Relation-first evidence memory}
\label{alg:appendix_relation_memory}
\begin{algorithmic}[1]
\REQUIRE Rollout result $R_t$, previous memory $M_{t-1}$
\STATE $M_t\leftarrow M_{t-1}$
\FORALL{record $v$ observed in $R_t$}
    \STATE Update average risk score and visit count for $v$
    \STATE Store why $v$ is risky in relation to visible predecessors or successors
\ENDFOR
\FORALL{reported relation $(u,v)$ in $R_t$}
    \STATE Update pair evidence that $(u,v)$ creates or supports a conflict
    \STATE Update second-endpoint evidence when $v$ becomes risky because of earlier evidence
    \STATE Update affected evidence for downstream records whose status changes
\ENDFOR
\IF{$R_t$ contains an \oc{} signal}
    \STATE Add its implicated endpoints and supporting relation to \oc{} memory
\ENDIF
\RETURN $M_t$
\end{algorithmic}
\end{algorithm}

\begin{algorithm}[H]
\small
\caption{Dynamic core compression}
\label{alg:appendix_dynamic_core}
\begin{algorithmic}[1]
\REQUIRE Previous core $H_{t-1}$, memory $M_t$, candidate cap $B(S)$
\STATE Form candidate set $C_t$ from prior core, local risky records, relation endpoints, affected records, and \oc{} support records
\FORALL{$v\in C_t$}
    \STATE Score $v$ by node risk, pair support, affected-region support, \oc{} support, and persistence
    \STATE Mark $v$ as protected if repeated evidence indicates a critical endpoint
\ENDFOR
\STATE Start $H_t$ with protected endpoints and strongest relation-supported records
\WHILE{$|H_t|<B(S)$ and useful candidates remain}
    \STATE Add the candidate with largest marginal support for explaining the risk
\ENDWHILE
\WHILE{$|H_t|>B(S)$}
    \STATE Remove the weakest non-protected candidate
\ENDWHILE
\RETURN $H_t$
\end{algorithmic}
\end{algorithm}

The core is therefore not a monotone union of all suspicious records.  Later rollouts can replace weak records with stronger evidence, although the update cannot drop repeatedly supported critical endpoints only to obtain a smaller subgraph.

\FloatBarrier
\section{Formula Details for Evidence Updates}
\label{sec:appendix_formula_details}

Implementation-level equations summarized in Section~\ref{sec:method} are expanded here.  The main text retains the search template plus two scoring components, the priority score and core objective.  Remaining details were fixed before evaluation and left untuned on test cases.

\paragraph{Node and relation updates.}
For every record observed in rollout $t$, node risk evidence is updated by a visit-counted running average:
\begin{equation}
    \begin{aligned}
    \bar{q}_{t+1}(v)=
    \frac{n_t(v)\bar{q}_t(v)+y_t(v)}
    {n_t(v)+1},\\
    n_{t+1}(v)=n_t(v)+1.
    \end{aligned}
    \label{eq:node-update}
\end{equation}
Directed relation evidence is updated by exponential smoothing,
\begin{equation}
    m_{t+1}(u,v)=(1-\rho)m_t(u,v)+\rho z_t(u,v),
    \label{eq:relation-update}
\end{equation}
where $z_t(u,v)$ is the local evaluator's evidence that relation $(u,v)$ creates or supports a repair-relevant conflict.

\paragraph{Relation priors and critical pairs.}
Equation~\ref{eq:relation-prior} is shared by node and edge priors.  For nodes we set $w=\theta$ and use node features such as repair-entry support, local-subgraph appearance, pair-endpoint support, and \oc{} endpoint evidence.  For edges we set $w=\beta$ and use edge features such as semantic incompatibility, repeated pair support, reverse-edge evidence, and accumulated conflict evidence.  A separate bounded ledger records persistent critical-pair hypotheses:
\begin{equation}
    \begin{aligned}
    \mathcal{N}(v)&=\{u:(u,v)\in E \text{ or } (v,u)\in E\},\\
    p_t(v)&=\max_{u\in\mathcal{N}(v)} \kappa_t(u,v),\\
    \kappa_t(u,v)&=\min(1,\omega^\top\chi_t(u,v)).
    \end{aligned}
    \label{eq:critical-pair}
\end{equation}
Here $\chi_t(u,v)$ contains pair-level evidence such as semantic incompatibility, severity, same-path support, no-fallback language, repeated support, and explicit conflict reports.

\paragraph{Online conflict signal.}
Let $c_t=\max_{(u,v)\in E(W_t)} z_t(u,v)$ be the strongest relation-conflict score in the current window.  With previous rollout scores as the online reference set,
\begin{equation}
    \tau_t = Q_{1-\alpha_{\mathrm{oc}}}\left(\{c_i: i<t\}\right),
    \quad
    \operatorname{OC}_t=\mathbf{1}[c_t>\tau_t].
    \label{eq:oc-rule}
\end{equation}
The implicated relation endpoints are added to \oc{} memory and receive the bounded bonus in Eq.~\ref{eq:priority}.

\paragraph{Core candidate pool.}
The candidate pool used in Eq.~\ref{eq:core} is
\begin{equation}
    C_t = h_t \cup H_{t-1} \cup O_t \cup K_t,
    \label{eq:candidate-pool}
\end{equation}
where $h_t$ is the local risk subgraph, $O_t$ contains \oc{} endpoints, and $K_t$ contains top-scoring relation endpoints together with their local neighbors.

\section{Fixed Search-Prior Coefficients}
\label{app:priors}

The numeric values in Table~\ref{tab:prior_coefficients} are fixed search-prior coefficients, not learned parameters.  They were selected using pilot checks before the locked evaluation.  Domain experts were used only as a sanity check on the ordinal reliability ordering among evidence types, not to tune coefficients on reported test cases.  Prior MCTS work motivates the structure, including UCB/UCT exploration \citep{kocsis2006bandit}, visit-decayed heuristic priors through progressive bias \citep{chaslot2008progressive,browne2012survey}, AMAF/RAVE-style evidence reuse \citep{gelly2007combining}, and graph/transposition-aware value reuse \citep{childs2008transpositions,czech2021monte}.  The coefficients instantiate these ideas for structural document risk.

\begin{table*}[t]
\centering
\small
\begin{tabular}{p{.20\textwidth}p{.10\textwidth}p{.62\textwidth}}
\toprule
Parameter & Value & Role and rationale \\
\midrule
\multicolumn{3}{l}{\textit{Node-feature weights }$\theta$ for Eq.~\ref{eq:relation-prior}}\\
$\theta_{\mathrm{repair}}$ & 1.10 & LLM names the record as a repair entry. This is closest to the final remediation target.\\
$\theta_{\mathrm{local}}$ & 0.90 & The record appears in local risk subgraphs; useful but noisier than explicit repair-entry evidence.\\
$\theta_{\mathrm{pair}}$ & 0.85 & The record is an endpoint of a supported conflict pair; structural risks often require two endpoints.\\
$\theta_{\mathrm{node}}$ & 0.70 & AMAF/RAVE-style node evidence reused across rollouts; important but indirect.\\
$\theta_{\mathrm{path}}$ & 0.75 & Path-fragment evidence that the record lies on repeatedly implicated local traversals.\\
$\theta_{\mathrm{conflict}}$ & 0.60 & Single-window conflict reports; useful but downweighted because local LLM judgments can be noisy.\\
$\theta_{\mathrm{OC}}$ & 2.00 & Effective ordered-conflict endpoint; high-confidence routing signal, later normalized and visit-decayed.\\
\midrule
\multicolumn{3}{l}{\textit{Edge-feature weights }$\beta$ for Eq.~\ref{eq:relation-prior}}\\
$\beta_{\mathrm{incompat}}$ & 1.55 & Semantic incompatibility is the strongest relation-level evidence for structural contradiction.\\
$\beta_{\mathrm{severity}}$ & 1.35 & Severity increases priority, but high severity alone is not sufficient without structural support.\\
$\beta_{\mathrm{explicit}}$ & 0.90 & Explicit conflict wording helps, but is downweighted relative to persistent structural evidence.\\
$\beta_{\mathrm{pair}}$ & 0.70 & Repeated pair support indicates relation-level consistency across windows.\\
$\beta_{\mathrm{edge}}$ & 0.65 & AMAF edge evidence reused from related traversals.\\
$\beta_{\mathrm{reverse}}$ & 0.35 & Reverse-edge support is weak auxiliary evidence.\\
$\beta_{\mathrm{conflict}}$ & 0.55 & The edge's own accumulated conflict estimate, kept below semantic incompatibility and severity.\\
\midrule
\multicolumn{3}{l}{\textit{Critical-pair prior }$\omega$ in Eq.~\ref{eq:critical-pair}}\\
$\omega_{\mathrm{sem}}$ & 0.30 & Semantic incompatibility dominates because it directly captures mutual exclusion.\\
$\omega_{\mathrm{sev}}$ & 0.25 & Severity is second: critical risks deserve revisits, but severity alone can be diffuse.\\
$\omega_{\mathrm{path}}$ & 0.18 & Same-path evidence strengthens the claim that two endpoints jointly define one repair region.\\
$\omega_{\mathrm{nofallback}}$ & 0.12 & No-fallback language makes contradictions harder to resolve locally.\\
$\omega_{\mathrm{support}}$ & 0.10 & Repeated support stabilizes initially uncertain pair hypotheses.\\
$\omega_{\mathrm{explicit}}$ & 0.05 & Explicit conflict reports are useful but lowest-weighted to avoid rewarding one noisy sentence.\\
\bottomrule
\end{tabular}
\caption{Fixed search-prior coefficients used by \samcgs{}.  Higher values reflect more reliable or more directly repair-relevant evidence.  Values were frozen before the locked evaluation and are not learned from the reported test cases.}
\label{tab:prior_coefficients}
\end{table*}

\FloatBarrier
\section{Shared Risk Rubric Prompt}
\label{sec:appendix_shared_rubric}

Reviewer fairness depends on separating the risk definition from the search mechanism.  The following excerpt gives the shared evaluator rubric used by \revnewreplace{both full-\scc{} \naive{} attempts and the \samcgs{} local-window evaluator}{one-shot prompting and the local-window evaluators used by LEA and \samcgs{}}.

\begin{quote}
\small
You are given records from a directed graph of interdependent documents.  Identify records or record sets that are structurally inconsistent, mutually incompatible, underspecified, or high-impact if left unresolved.  Use only visible record text and visible directed relations.  Do not assume a domain-specific defect type.  Return evidence-based judgments with node ids and concise reasons.
\end{quote}

\paragraph{Shared semantics.}
\revredreplace{\naive{} and \samcgs{} use this same risk rubric.  The methods differ in input unit and search procedure.  Each \naive{} attempt reads the complete \scc{} and emits a global ranking plus a risk subgraph; the main baseline aggregates ten such attempts with oracle-risk Top-3 selection.  \samcgs{} reads local windows over multiple rollouts and aggregates evidence before emitting a final core.  The comparison holds the risk definition fixed and varies the computation.}{One-shot prompting, LEA, and \samcgs{} share this risk rubric. One-shot produces a global ranking and subgraph in one call; LEA aggregates fixed local evaluations; \samcgs{} adaptively acquires local evidence to construct its final core.}

\section{Prompt Templates}
\label{sec:appendix_prompts}

Complete prompts include serialization details and JSON schemas.  The abstract templates reported here are sufficient to reproduce the evaluation interface without depending on one implementation file.

\paragraph{\revnewreplace{\naive{} full-\scc{} direct-subgraph prompt.}{One-shot prompting template.}}
\begin{quote}
\small
Input: the complete node list for one strongly connected component, the text of every node, and all directed edges among the nodes.  Apply the shared risk rubric to the full component.  Return valid JSON with \texttt{global\_ranking}, \texttt{risk\_subgraph\_nodes}, \texttt{risk\_factor\_distribution}, and \texttt{explanation}.  Rankings and subgraph selections should be grounded in node IDs and visible edges.
\end{quote}

\paragraph{\samcgs{} local-window prompt.}
\begin{quote}
\small
Input: a local window from one strongly connected component, including visible node text, visible directed edges, and a concise prior evidence summary.  Apply the same shared risk rubric to the local window.  Return valid JSON with \texttt{local\_node\_scores}, \texttt{relation\_evidence}, \texttt{local\_risk\_subgraph\_nodes}, \texttt{repair\_entry\_nodes}, optional \texttt{oc\_candidate}, and a concise explanation grounded in exact node IDs.
\end{quote}

\paragraph{Output contrast.}
\revredreplace{\naive{} requires each generation to rank the entire \scc{} and produce a final subgraph.  In contrast, \samcgs{} asks only for local evidence, leaving final subgraph construction to algorithmic memory and compression.}{One-shot prompting ranks the entire component and returns a final subgraph in one generation. LEA and \samcgs{} share a local evaluator interface and aggregate its outputs algorithmically. LEA uses fixed windows without adaptive search or memory; Appendix~\ref{app:revision_lea} specifies its window construction and aggregation.}

\section{Dataset and Injection Details}
\label{sec:appendix_datasets}

\begin{revredblock}
\paragraph{Formatting interface.}
Record formats vary by domain: Debian records contain package fields, BGB records are statutory sections, SEC EX-21 records are subsidiary disclosures, and CUAD records are contract clauses. Lightweight domain formatters
\begin{equation}
    \mathcal{A}_d=(\operatorname{serialize}_d,\operatorname{edges}_d,\operatorname{normalize}_d)
\end{equation}
serialize records, expose typed directed edges, and normalize model outputs without specifying a target risk type or revealing injected endpoints. All three methods apply the same domain-neutral rubric to visible node text and directed edges, identifying structural inconsistency, mutual incompatibility, underspecification, and high-impact risks. They differ in how they acquire and aggregate evidence: global generation, fixed local aggregation, or adaptive \scc{}-aware search.
\end{revredblock}


Table~\ref{tab:dataset_processing} summarizes the four included domains and their conversion into directed record graphs.  Domain coverage spans several forms of authority and structure, including technical dependency metadata, public corporate disclosures, official statutory references, and expert-labeled contract clauses.

\begin{table*}[t]
\centering
\scriptsize
\setlength{\tabcolsep}{3pt}
\begin{tabular}{p{.12\textwidth}p{.20\textwidth}p{.30\textwidth}p{.30\textwidth}}
\toprule
Domain & Why selected & Graph construction & Role in evaluation \\
\midrule
Debian & Technical dependency graph with explicit package relationships. & Package metadata entries become records; dependency, replacement, conflict, and virtual-provider fields become directed edges. & \revredreplace{Tests whether the method generalizes beyond legal and financial language.}{Evaluates software migration and dependency constraints beyond legal and financial text.} \\
SEC EX-21 & Public company subsidiary disclosures with entity-control structure. & Subsidiary names, jurisdictions, consolidation fields, ownership hints, and control relations become entity records and directed edges. & \revredreplace{Tests public regulatory disclosure reasoning and ownership-chain consistency.}{Evaluates ownership and consolidation reasoning from public financial and legal disclosures.} \\
BGB & Official German Civil Code cross-references. & Statutory sections become records; references, obligations, conditions, exceptions, and remedies become directed edges. & \revredreplace{Provides a comparatively clean legal-rule stress test with authoritative source text.}{Provides a relatively clean setting for statutory cross-reference cycles in authoritative legal text.} \\
CUAD & Expert-labeled commercial contract dataset. & Clauses and definitions become records; definition use, cross-reference, amendment, condition, termination, and obligation relations become edges. & \revredreplace{Provides long, noisy contract components where evidence can be diffuse.}{Tests long contract components with ambiguity and diffuse evidence.} \\
\bottomrule
\end{tabular}
\caption{\revnewreplace{Dataset construction and processing.}{Domain coverage, dataset construction, and evaluation roles.} All domains are represented as directed record graphs before \scc{} localization.}
\label{tab:dataset_processing}
\label{tab:domains}
\end{table*}

\paragraph{SCC construction and filtering.}
For each domain, we construct a directed record graph and enumerate non-trivial \scc{}s using Tarjan's algorithm.  Candidate components are filtered by size and record availability and then locked before model evaluation.  The main experiment uses the same locked \scc{} blocks \revnewreplace{for \naive{} and \samcgs{} under the same model and profile.}{for One-shot, LEA, and \samcgs{} under the same model and profile.}

\paragraph{Injection constraints.}
The critical benchmark modifies only existing records in the selected \scc{}.  No new nodes are added.  The injected root record provides the endpoint used for \rootthree{} evaluation.  A distant witness record supplies the evidence that makes the root-side statement structurally incompatible under the graph context.  Affected records are downstream or neighboring records whose interpretation changes under the root--witness conflict.  Subgraph metrics count root, witness, and affected records as repair-relevant endpoints when applicable. \revred{Table~\ref{tab:injection_templates} lists the four injection templates.}

\begin{table*}[t]
\centering
\scriptsize
\setlength{\tabcolsep}{3pt}
\begin{tabular}{p{.16\textwidth}p{.24\textwidth}p{.26\textwidth}p{.24\textwidth}}
\toprule
Template & Root-side edit & Witness / affected edit & Intended risk \\
\midrule
\texttt{direct\_mutex} & Adds a direct obligation, ownership, compatibility, or condition statement expected to hold. & Adds a distant statement that requires the opposite state under the same dependency chain. & Two records are individually plausible but mutually incompatible when the component context is followed. \\
\texttt{handoff\_invariant} & States that an invariant is expected to be preserved before control, migration, or responsibility is handed off. & States that the downstream record can proceed only after the invariant has been changed or removed. & The handoff cannot be completed because the precondition and downstream entry condition disagree. \\
\texttt{temporal\_gate} & Specifies a deadline, waiting period, migration window, vesting condition, or trigger order. & Requires a conflicting temporal gate, such as earlier execution, delayed activation, or a missing prerequisite. & The same chain imposes incompatible timing constraints. \\
\texttt{condition\_trigger} & Makes a right, remedy, interface, or ownership statement conditional on a trigger. & Makes the same trigger impossible, optional, or reversed for a dependent record. & A dependent record relies on a condition whose truth value is changed elsewhere in the component. \\
\bottomrule
\end{tabular}
\caption{Critical structural injection templates.  All templates edit existing records only and avoid explicit defect labels.}
\label{tab:injection_templates}
\end{table*}

\paragraph{Why \riskall{} matters.}
\riskany{} is useful when a repair process can sometimes start from one good entry point.  \riskall{} is stricter and central for critical cases, where fixing a structural contradiction often requires both sides of the incompatibility.  A subgraph that includes only the obvious-risk paragraph but omits the distant condition may be too small to support a reliable repair.

\FloatBarrier
\section{Additional Experimental Details}
\label{sec:appendix_experiments}

\begin{revnewblock}
\paragraph{Locked run configuration.}
\revred{All five models are evaluated on the same 80 injected variants from 20 source \scc{}s. One-shot, LEA, and \samcgs{} use the shared risk rubric and fixed method configurations specified in Appendix~\ref{app:revision_lea}. Table~\ref{tab:appendix_breakdowns} groups results by model, domain, and \scc{} size.}
\end{revnewblock}

\input{revision_retired_controls}

\begin{revnewblock}
\input{../../experiments/main_experiment/paper_assets/tables/five_model_appendix_breakdowns}
\end{revnewblock}

% October 1: both size views and complete size analysis are retained in Appendix J.1.
\begin{revredblock}
\subsection{Performance by SCC size}
\label{app:size_breakdowns}

\samcgs{} improves \riskall{} over LEA in all three \scc{}-size groups for GPT-4o and Qwen2.5-72B (Figure~\ref{fig:main_metrics}). Their largest gains occur in medium components: 37.50 and 62.50 percentage points, respectively. For Llama 3.1-8B-Instruct, \rootthree{} also improves in every size group. Absolute results appear in Figure~\ref{fig:absolute_size_breakdowns} and Table~\ref{tab:appendix_breakdowns}.

\begin{figure*}[t]
    \centering
    \includegraphics[width=\textwidth]{fig02_scale_gain_heatmaps.pdf}
    \caption{\revredreplace{Main critical-\scc{} results on 320 model-case pairs per method. Panel A compares \samcgs{} with oracle-risk \naive{}, an oracle-risk Top-3 selection over ten full-\scc{} attempts; Panel B gives a compact \scc{}-size breakdown of endpoint retention and compression.}{\samcgs{} gains over LEA by \scc{} size: (a) \rootthree{} and (b) \riskall{}. Cells show percentage-point differences for each model and size group (36 small, 24 medium, and 20 large cases per model). Blue indicates gains and orange indicates losses on a shared zero-centered scale. Figure~\ref{fig:absolute_size_breakdowns} reports absolute performance for all three methods.}}
    \label{fig:main_metrics}
\end{figure*}

Figure~\ref{fig:absolute_size_breakdowns} uses the same size groups and case counts as Figure~\ref{fig:main_metrics}, reporting \rootthree{} and \riskall{} for all three methods and five models.

\begin{figure*}[t]
    \centering
    \includegraphics[width=\textwidth]{fig02_main_results_five_models.pdf}
    \caption{Performance by \scc{} size. Columns correspond to models and rows to (a) \rootthree{} and (b) \riskall{}. All three methods use the same small ($\leq12$ nodes), medium (14--20), and large ($\geq24$) case groups. Gray circles, orange squares, and blue triangles denote One-shot, LEA, and \samcgs{}, respectively.}
    \label{fig:absolute_size_breakdowns}
\end{figure*}
\end{revredblock}

% M4: retain the complete budget-prefix plot and numeric analysis in Appendix J.
\begin{revnewblock}
\subsection{Budget-prefix convergence}
\label{app:budget_prefix}

\revred{Most gains in evidence completeness occur within the first 20 rollouts (Figure~\ref{fig:convergence}): \riskall{} rises from 25.00\% at budget 5 to 76.25\% at budget 20. Subsequent rollouts refine the core, reaching 75.00\% \riskall{} and 51.26\% compression at budget 60, matching the GPT-4o results in Table~\ref{tab:matched_main}.}

\begin{figure*}[t]
    \centering
    \includegraphics[width=\textwidth]{\revasset{fig04_budget_prefix_convergence.pdf}{fig04_budget_prefix_three_methods.pdf}}
    \caption{\revredreplace{Budget-prefix convergence: localization and evidence coverage improve with rollout budget while the retained core stays compact.}{GPT-4o evidence retention and compression across 60 rollout prefixes on the same 80 cases. The solid curve represents \samcgs{}; horizontal lines indicate the final One-shot and LEA scores. Each prefix is reconstructed from stored search evidence using the final-core selection rule. Results at budget 60 match the final outputs.}}
    \label{fig:convergence}
\end{figure*}
\end{revnewblock}

\FloatBarrier
\input{revision_experiments}

\end{document}
